\clearpage\SourceLeaf{190}
\begin{center}
{\Large\textls{ANZEIGE.}}
\end{center}

\bigskip
\noindent\hrulefill

\begin{center}
Göttingische gelehrte Anzeigen. 1822 Oktober 28.
\end{center}

\noindent\hrulefill
\bigskip

\begin{center}
\textls{Marburg.}
\end{center}

\emph{Theorie der Parallelen, von} \textsc{Carl Reinhard Müller}, \emph{Doktor der Philosophie, außerordentlichem Professor der Mathematik} und\,so\,weiter. 1822. 40 Seiten im Quartformat.

Der Rezensent hat bereits vor sechs Jahren in diesen Blättern seine Überzeugung ausgesprochen, dass alle bisherigen Versuche, die Theorie der Parallellinien streng zu beweisen oder die Lücke in der euklidischen Geometrie auszufüllen, uns diesem Ziel nicht nähergebracht haben. Er kann nicht anders, als dieses Urteil auch auf alle späteren ihm bekannt gewordenen Versuche auszudehnen. Dennoch bleiben manche solcher Versuche, obwohl sie ihren eigentlichen Hauptzweck verfehlen, wegen des darin bewiesenen Scharfsinns für Freunde der Geometrie lesenswert. Der Rezensent glaubt, in dieser Hinsicht die vorliegende, anlässlich einer Schulprüfung veröffentlichte kleine Schrift besonders hervorheben zu müssen. Den ganzen einfallsreichen Gedankengang des Verfassers hier ausführlich darzulegen, wäre für unsere Blätter zu weitläufig und auch überflüssig, da die Schrift selbst gelesen zu werden verdient. Sie hat aber wie alle übrigen Versuche ihre schwache Stelle. Diese herauszuheben ist der Zweck dieser Anzeige.

Wir finden diese schwache Stelle auf S.\,15 im Beweis des Lehrsatzes des 15. Artikels. Dieser Lehrsatz ist der eigentliche Nerv der ganzen Theorie, die fällt, sobald er sich nicht streng beweisen lässt. Wir führen daher zunächst diesen Lehrsatz an. Die zugehörige Figur kann jeder leicht selbst zeichnen.

Wenn jeder Winkel an der Grundlinie $ON$ eines gleichschenkligen Dreiecks größer ist als der Winkel an der Spitze $A$ und man in $O$ an die Seite $OA$
